\documentclass{article}
\usepackage[T1]{fontenc}
\usepackage{amsmath, amssymb, amsthm}
\usepackage[mathscr]{euscript}
\usepackage{fullpage}
\usepackage[utf8]{inputenc}
\usepackage[french]{babel}
\usepackage{geometry}
\geometry{a4paper, margin=2.5cm}

\newcommand{\Hom}{\mathrm{H}}
\newcommand{\ZZ}{\mathbb{Z}}

\newtheorem{definition}{Définition}
\newtheorem{proposition}{Proposition}
\newtheorem{theorem}{Théorème}

\DeclareMathOperator{\sign}{sign}
\DeclareMathOperator{\Bl}{Bl}
\newcommand{\CP}{\mathbb{C}\mathbb{P}}
\newcommand{\RP}{\mathbb{R}\mathbb{P}}
\newcommand{\Z}{\mathbb{Z}}
\newcommand{\C}{\mathbb{C}}
\newcommand{\R}{\mathbb{R}}
\newcommand{\Htwo}{H_2}
\newcommand{\HtwoR}{H_2(-;\R)}

\begin{document}

\title{Forme d'intersection, signature et éclatement}
\author{}
\date{}
\maketitle

\section*{Cas \(k = 2\) : homologie de degr\'e 2}

Soit \(M_1^4\) et \(M_2^4\) deux vari\'et\'es connexes, orientables, de dimension 4. On consid\`ere leur somme connexe \(M = M_1 \# M_2\).

On pose \(U = M_1 \setminus \{\text{int\'erieur d'une petite boule}\}\) et \(V = M_2 \setminus \{\text{int\'erieur d'une petite boule}\}\), l\'eg\`erement agrandis pour former un recouvrement de \(M\). L'intersection \(U \cap V\) est hom\'eomorphe \`a \(\mathbb{S}^3 \times (-\epsilon, \epsilon)\), donc homotopiquement \'equivalente \`a \(\mathbb{S}^3\).

La suite exacte de Mayer-Vietoris en homologie r\'eduite donne :
\[
\cdots \to \Hom_2(\mathbb{S}^3) \to \Hom_2(U) \oplus \Hom_2(V) \to \Hom_2(M) \to \Hom_1(\mathbb{S}^3) \to \cdots
\]

On conna\^it l'homologie de la sph\`ere \(\mathbb{S}^3\) :
\[
\Hom_1(\mathbb{S}^3) = 0 \quad \text{et} \quad \Hom_2(\mathbb{S}^3) = 0.
\]

De plus, \(U\) se r\'etracte par d\'eformation sur \(M_1\) priv\'e d'un point, donc \(\Hom_2(U) \cong \Hom_2(M_1)\). De m\^eme, \(\Hom_2(V) \cong \Hom_2(M_2)\). (On pourrait aussi utiliser le théorème d'excision)

La suite exacte devient alors :
\[
0 \to \Hom_2(M_1) \oplus \Hom_2(M_2) \to \Hom_2(M) \to 0.
\]

Cette suite \'etant exacte, la fl\`eche du milieu est un isomorphisme. On obtient donc :
\[
\boxed{\Hom_2(M_1 \# M_2) \cong \Hom_2(M_1) \oplus \Hom_2(M_2)}.
\]

\section{Forme d'intersection}

Soit $M$ une variété lisse, fermée, orientable de dimension $4$. La forme d'intersection est un outil fondamental qui encode l'intersection des cycles de dimension $2$.

\begin{definition}
Soit $M$ une $4$-variété lisse fermée orientée. La \textbf{forme d'intersection}
\[
Q_M : H_2(M;\Z) \times H_2(M;\Z) \longrightarrow \Z
\]
est définie par
\[
Q_M(\alpha,\beta) = \langle \alpha \cup \beta, [M] \rangle,
\]
où $\alpha,\beta \in H^2(M;\Z)$ sont les classes de Poincaré duales de $\alpha$ et $\beta$ dans la cohomologie. 
\end{definition}

On préfère parfois une description alternative. Soit $\alpha, \beta \in H_2(M;\Z)$. Choisissons des surfaces lisses orientées $A,B \subset M$ représentant $\alpha$ et $\beta$ et les mettant en position générale. Alors $A$ et $B$ s'intersectent transversalement en un nombre fini de points. En orientant $M$ et les surfaces, chaque point d'intersection reçoit un signe $\pm 1$. La somme de ces signes est $Q_M(\alpha,\beta)$, indépendante des choix.

\begin{definition}[Version alternative]
Soit $\alpha,\beta \in H_2(M;\Z)$. Choisissons des cycles lisses orientés $a,b$ représentant $\alpha,\beta$ transverses. Alors
\[
Q_M(\alpha,\beta) = \sum_{p \in a \cap b} \varepsilon(p),
\]
où $\varepsilon(p) = \pm 1$ est le signe d'intersection déterminé par les orientations.
\end{definition}

La forme $Q_M$ est symétrique ($Q_M(\alpha,\beta)=Q_M(\beta,\alpha)$) et unimodulaire (car $M$ est fermée orientée). En coefficients réels, $Q_M$ est une forme bilinéaire symétrique sur $H_2(M;\R)$.

\begin{definition}
La \textbf{signature} $\sigma(M)$ de $M$ est la signature de la forme $Q_M$ sur $H_2(M;\R)$, i.e. $\sigma(M) = b_2^+(M) - b_2^-(M)$, où $b_2^+(M)$ (resp. $b_2^-(M)$) est la dimension d'un sous-espace maximal sur lequel $Q_M$ est définie positive (resp. négative).
\end{definition}

\section{Exemple : $\mathbb{CP}^2$ et $\mathbb{CP}^2 \# n\overline{\mathbb{CP}^2}$}

\begin{itemize}
    \item Pour $\CP^2$ (orientation complexe standard), $H_2(\CP^2;\Z) \cong \Z$ engendré par la classe de la droite projective $\ell$. On a $\ell \cdot \ell = 1$. Donc $Q_{\CP^2} = (1)$ et $\sigma(\CP^2)=1$.
    \item Pour $\overline{\CP}^2$ (même variété avec orientation opposée), $Q_{\overline{\CP}^2} = (-1)$ et $\sigma = -1$.
\end{itemize}

Soit $X_n = \CP^2 \# n\overline{\CP}^2$, la somme connexe de $\CP^2$ avec $n$ copies de $\overline{\CP}^2$. Par additivité de la forme d'intersection sous somme connexe en dimension $4$ (à condition de prendre en compte l'orientation), on a
\[
H_2(X_n;\Z) \cong \Z^{n+1},
\]
et la forme d'intersection est la somme orthogonale
\[
Q_{X_n} = \langle 1 \rangle \oplus \underbrace{\langle -1 \rangle \oplus \cdots \oplus \langle -1 \rangle}_{n \text{ fois}}.
\]
En termes de matrice,
\[
Q_{X_n} = \begin{pmatrix}
1 & & & \\
& -1 & & \\
& & \ddots & \\
& & & -1
\end{pmatrix}.
\]
Sa signature est donc
\[
\sigma(X_n) = 1 - n.
\]

\section{Éclatement (blow-up) de $\mathbb{C}^n$ en $0$}

L'éclatement est une opération locale qui remplace un point par un diviseur exceptionnel.

\begin{definition}
Soit $U = \C^n$ avec l'origine $0$. Le \textbf{blow-up} de $\C^n$ en $0$ est la variété
\[
\Bl_0 \C^n = \left\{ (x,\ell) \in \C^n \times \CP^{n-1} \mid x \in \ell \right\}.
\]
La projection $\pi : \Bl_0 \C^n \to \C^n$ est donnée par $(x,\ell) \mapsto x$.
\end{definition}


En dimension $4$, l'éclatement d'une $4$-variété en un point correspond à la somme connexe avec $\overline{\CP}^2$. Plus précisément, si $M$ est une $4$-variété lisse orientée et $p \in M$, alors
\[
M \# \overline{\CP}^2 \cong \Bl_p M.
\]
Ainsi, la forme d'intersection se modifie par addition d'une copie de $\langle -1 \rangle$ :
\[
Q_{\Bl_p M} \cong Q_M \oplus \langle -1 \rangle.
\]
Cette relation explique pourquoi $\CP^2 \# n\overline{\CP}^2$ est obtenu en éclatant $n$ fois $\CP^2$ en des points génériques.

\end{document}